\documentclass[../main.tex]{subfiles}

\begin{document}
%% begin abstract format
\makeatletter
\renewenvironment{abstract}{%
    \if@twocolumn
      \section*{Resumen \\}%
    \else %% <- here I've removed \small
    \begin{flushright}
        {\filleft\Huge\bfseries\fontsize{48pt}{12}\selectfont Resumen\vspace{\z@}}%  %% <- here I've added the format
        \end{flushright}
      \quotation
    \fi}
    {\if@twocolumn\else\endquotation\fi}
\makeatother
%% end abstract format
%% begin abstract format
\makeatletter
\renewenvironment{abstract}{%
    \if@twocolumn
      \section*{Resumen \\}%
    \else %% <- here I've removed \small
    \begin{flushright}
        {\filleft\Huge\bfseries\fontsize{48pt}{12}\selectfont Resumen\vspace{\z@}}%  %% <- here I've added the format
        \end{flushright}
      \quotation
    \fi}
    {\if@twocolumn\else\endquotation\fi}
\makeatother
%% end abstract format

\begin{abstract}

\vspace{30pt}

Este Trabajo de Fin de Grado aborda el diseño, implementación y validación experimental de un prototipo de casco inteligente orientado a la monitorización de impactos craneales en deportes de combate. La motivación del proyecto se fundamenta en la creciente preocupación por las lesiones cerebrales traumáticas asociadas a impactos repetidos en la cabeza y en la necesidad de herramientas tecnológicas que permitan cuantificar objetivamente la exposición al riesgo.

El sistema desarrollado integra una plataforma embebida portable con comunicación Bluetooth Low Energy (BLE) y una unidad de medición inercial (IMU) capaz de registrar aceleraciones lineales y velocidades angulares. Mediante una arquitectura software basada en adquisición por eventos, el dispositivo detecta impactos, estima la aceleración rotacional máxima y clasifica cada evento en rangos de severidad biomecánica definidos a partir de la literatura científica.

El prototipo se integró en un casco de protección mediante un encapsulado diseñado en CAD e impreso en 3D, y fue evaluado en sesiones reales de entrenamiento con varios deportistas. Los resultados muestran una detección consistente bajo condiciones adecuadas de fijación mecánica y evidencian la importancia del acoplamiento sensor--casco para garantizar la fiabilidad de la medida.

El sistema no constituye una herramienta diagnóstica clínica, sino una solución orientada a la monitorización preventiva y a la gestión prudente de la exposición acumulada a impactos craneales en el ámbito deportivo.

\vspace{12pt}

\bfseries{\large{Palabras clave:}} lesión cerebral traumática, aceleración rotacional, monitorización de impactos craneales, dispositivos wearables, sistemas embebidos.

\end{abstract}


\end{document}


